\section{The actual cubic characters with a rational cubic seven-factor}
\label{sec:mixed-cubic-seven-cubic-strict}

Put $K=\mathbb Q(\omega)$, $t=\sqrt[3]{269}$, $L=K(t)$,
and $R=\mathbb Q(\zeta_7)^+$.  Let $\chi$ be either faithful
character of $\operatorname{Gal}(L/K)$, and let $\nu$ be
either faithful character of $\operatorname{Gal}(R/\mathbb Q)$.
In the actual conductor-seven notation, $\nu=\xi_7^2$ for
a choice of primitive order-six character $\xi_7$ modulo seven.
Both powers of each character will be included, so these
choices do not affect the stated set of labels.
Write
\[
              F=LR,\qquad B=\mathbb Q(t)R.
 \tag{MC7.1}
\]
For a prime $\ell>3$ with $\ell\ne7,269$, define
$c(\ell)=1$ exactly when $\ell\equiv1\pmod3$ and
$269^{(\ell-1)/3}\not\equiv1\pmod\ell$; in every other
case put $c(\ell)=0$.  Define $s(\ell)=0$ when
$\ell\equiv1$ or $-1\pmod7$, and $s(\ell)=1$ otherwise.
Use the explicit prime set
\[
 \begin{split}
 \mathscr P_7=\{\ell\text{ prime}:\ &\ell>3,\quad\ell\ne7,269,\\
                                    &c(\ell)+s(\ell)=1\}.
 \end{split}
 \tag{MC7.2}
\]
Thus precisely one of the local cubic Kummer action and the
local rational cubic seven-action will be nontrivial.
This includes the split noncube primes congruent to $\pm1$
modulo seven, and also the complementary case where the
Kummer action is trivial and the rational cubic action is not.

For $\ell\in\mathscr P_7$, let $O$ be the original ring of
integers of a finite extension $E/\mathbb Q_\ell$ containing
$\mu_3$, with uniformizer $\varpi$ and residue field $k$.
Let $\kappa:G_K\to1+\mathfrak m_O$ be continuous and
unramified away $\ell$.  For $j,a\in\{1,2\}$ and $e\geq1$
divisible by $\ell-1$, set
\[
            D=(E/O)(\chi^j\nu^a\kappa^{-1}c_\ell^{-e}).
 \tag{MC7.3}
\]
Use the full strict kernel in $H^1(G_K,D)$: zero restriction
to $H^1(K_v,D)$ at every $v\mid\ell$, and to $H^1(I_v,D)$
at every other finite place.  Use the same conditions for
the original finite module $D[\varpi]$.

\begin{theorem}[AC-G209: original-coefficient strict zero with a cubic seven-factor]
\label{AC-G209}
For all the choices just stated,
\[
             \mathcal S_\ell(K,D[\varpi])=0,
             \qquad\mathcal S_\ell(K,D)=0.
 \tag{MC7.4}
\]
In particular, on the original growing host (FG.3) and original
coefficient orders (FC.1), all four actual labels
$\chi^j\xi_7^{2a}$, $j,a\in\{1,2\}$, have full strict zero
at these primes and weights.  Their two-power and auxiliary
$q,d$ factors are trivial.

For every integer $T\geq2$, define the actual triple set
\[
 \begin{split}
 \mathcal C_7(T)=\{(\ell,\chi^j\xi_7^{2a},e):\ &\ell\in\mathscr P_7,
   \ j,a\in\{1,2\},\\
 &1\leq e<T,\quad\ell-1\mid e\}.
 \end{split}
 \tag{MC7.5}
\]
Its exact cardinality is
\[
 |\mathcal C_7(T)|=
 4\sum_{\substack{\ell\in\mathscr P_7\\\ell\leq T}}
                 \left\lfloor\frac{T-1}{\ell-1}\right\rfloor.
 \tag{MC7.6}
\]
The primes $5,11,13$ belong to $\mathscr P_7$, giving the
explicit disjoint contribution
$4\lfloor(T-1)/4\rfloor+4\lfloor(T-1)/10\rfloor
+4\lfloor(T-1)/12\rfloor$.
For any retained set $J$ of original triples, the remainder
in (CP.3) obeys the exact equality
\[
                   \mathcal R_J=\mathcal R_{J\setminus\mathcal C_7(T)}.
 \tag{MC7.7}
\]
This equality concerns the full original $K$-strict summands
with nontrivial cubic factor.  It does not split positive
derived indices or assert that an individual such index is zero.
\end{theorem}

\begin{proof}
\proofstep{1. The exact degree-nine field and completed class quotient.}
The splitting field $L$ of $X^3-269$ has group $S_3$ over
$\mathbb Q$, with quadratic subfield $K$.
The field $R/\mathbb Q$ is cyclic of degree three.
Their intersection is Galois over $\mathbb Q$ and is either
$\mathbb Q$ or $R$.  The latter would give an order-three
quotient of $S_3$, which does not exist.  Hence
\[
 \operatorname{Gal}(F/\mathbb Q)=S_3\times C_3,
 \qquad\operatorname{Gal}(F/K)=C_3\times C_3.
 \tag{MC7.8}
\]
Let $\tau$ fix $t,R$ and invert $\omega$.  Since
$F=B(\omega)$ and this nontrivial automorphism fixes $B$,
one has $[F:B]=2$.  Thus $B=F^{\langle\tau\rangle}$
and $[B:\mathbb Q]=9$.

For $r=\zeta_7+\zeta_7^{-1}$ one has
$r^3+r^2-2r-1=0$.  The exact polynomial of $u=t+r$ is
\[
 \begin{split}
 Q(X)&=\operatorname{Res}_Y
          \bigl(Y^3+Y^2-2Y-1,(X-Y)^3-269\bigr)\\
 &=X^9+3X^8-3X^7-821X^6-1614X^5-3207X^4\\
 &\hspace{1cm}+212517X^3+209811X^2+1082181X-19172707.
 \end{split}
 \tag{MC7.9}
\]
Its nine roots are $t\omega^i+r_b$, where the $r_b$ are
the three conjugates of $r$ and $i=0,1,2$.
These roots are distinct.  If two roots with different
cubic indices were equal, then $t$ would belong to $KR$.
But $KR/\mathbb Q$ is abelian; since it contains $K$, such
an inclusion would put the nonabelian splitting field $L$
inside it, a contradiction.  With equal cubic indices the
conjugates $r_b$ must be equal.  Thus $u$ generates exactly
$B$, and (MC7.9) is its irreducible degree-nine polynomial.

The completed root-owned PARI computation uses this polynomial
in $B_0=\operatorname{bnfinit}(Q,1)$, followed by
$\operatorname{bnfcertify}(B_0,1)$.  The exact retained output is
\[
 \begin{array}{c|c}
 B_0.\mathrm{sign}&[3,3]\\
 B_0.\mathrm{no}&28\\
 B_0.\mathrm{cyc}&[14,2]\\
 \operatorname{bnfcertify}(B_0,1)&1.
 \end{array}
 \tag{MC7.10}
\]
The exact input, complete output, saved structure and execution
receipt are retained.  The process finished with exit code zero
at 09:33:07.387011 UTC on 23 November 2026, including the
successful completion marker.  The documented flag-one scope
is that the actual class group is a quotient of the computed
group.  Therefore
\[
 \begin{gathered}
 \operatorname{Cl}(B)\text{ is a quotient of }C_{14}\times C_2,\\
 \operatorname{Cl}(B)\{\ell\}=0
                         \quad(\ell>3,\ \ell\ne7).
 \end{gathered}
 \tag{MC7.11}
\]
No equality of actual class numbers or complete unit basis is
asserted.  The author inspects the completed root execution;
the documented software contract is an explicit input, not
a newly supplied proof of the entire class-group algorithm.

\proofstep{2. The targeted class action on the larger cover.}
Let $g$ generate the Kummer $C_3$ factor of
$\operatorname{Gal}(F/K)$, and let $h$ generate its rational
$R$ factor.  The field automorphism $\tau$ satisfies
\[
       \tau g\tau^{-1}=g^{-1},\qquad\tau h=h\tau.
 \tag{MC7.12}
\]
Put $C=\operatorname{Cl}(F)\{\ell\}$.
For the quadratic extension $F/B$, norm and extension of
ideal classes obey $N\circ j=2$ and $j\circ N=1+\tau$.
Since $\operatorname{Cl}(B)\{\ell\}=0$ and two is invertible
on $C$, one has $C^{\tau=1}=0$, equivalently $\tau=-1$ on
$C$.  Relation (MC7.12) then gives $g=g^{-1}$ as operators
on $C$.  With $g^3=1$ this forces $g=1$ on $C$.
Thus every class component with nontrivial Kummer character
is zero, regardless of its character on the rational $C_3$.
This does not claim that all of $C$ is zero.

\proofstep{3. Residual restriction and every local coefficient group.}
The original finite module is
\[
                     D[\varpi]=k(\overline\chi^{\,j}\overline\nu^{\,a}).
 \tag{MC7.13}
\]
Indeed $\kappa$ reduces to one, the reduction of $c_\ell^e$
is one, and reduction on $\mu_3$ is injective.
Each displayed residual character has order three, rather
than nine.  Its kernel inside $C_3\times C_3$ has order
three; the degree-nine cover $F/K$ simultaneously trivializes
all four characters and need not be their minimal splitting
field.

Restriction $H^1(G_K,D[\varpi])\to H^1(G_F,k)$ is injective,
since its inflation kernel is a cohomology group of the
order-nine group, killed by averaging in characteristic
$\ell>3$.  For a full strict residual class the restricted
homomorphism is unramified at every finite place of $F$ and
zero on the local groups above $\ell$.  As $F$ is totally
imaginary, it factors through $\operatorname{Cl}(F)$ and
hence through $C$.  Artin equivariance gives
\[
                  f(gx)=\overline\chi^{\,j}(g)f(x).
 \tag{MC7.14}
\]
The action of $g$ on $C$ is trivial by Step 2, while the
scalar $\overline\chi^{\,j}(g)$ is nontrivial.  Thus $f=0$,
proving the first assertion of (MC7.4).  The inverse convention
selects the other nontrivial cubic eigenvalue and gives the
same conclusion.

It remains to check every completion above the coefficient
prime.  For $v\mid\ell$, put $f_v=[K_v:\mathbb Q_\ell]$.
Since $\ell>3$, this is one at both places for split primes,
and two at the unique place for inert primes.
If $c(\ell)=1$, the prime is split and $X^3-269$ has
separable irreducible reduction at both places.  Thus
$\overline\chi^{\,j}(\operatorname{Frob}_v)$ has order three.
If $c(\ell)=0$ and the prime is split, then $269$ is a cube
modulo $\ell$, so a simple root lifts to $\mathbb Q_\ell$;
together with $\mu_3\subset K_v$ this splits the cubic.
If the prime is inert, cubing is bijective on
$\mathbb F_\ell^\times$, because $3\nmid\ell-1$.
Again a simple root lifts to $\mathbb Q_\ell$, and adjoining
$\mu_3\subset K_v$ supplies all three roots.  Hence
$\chi$ is trivial on $G_{K_v}$ in every case $c(\ell)=0$.

Frobenius in $R/\mathbb Q$ is the class of $\ell$ in
$(\mathbb Z/7\mathbb Z)^\times/\{1,-1\}$.
At a place of $K$ the rational character therefore has value
$\nu(\ell)^{f_v}$.  Its order is one if $s(\ell)=0$ and
three if $s(\ell)=1$, since $f_v=1$ or two is prime to three.
Exactly one of $c(\ell),s(\ell)$ is one.  Thus, at every
$v\mid\ell$ and for all four pairs $j,a$,
\[
 \overline\chi^{\,j}(\operatorname{Frob}_v)
       \overline\nu^{\,a}(\operatorname{Frob}_v)\ne1.
 \tag{MC7.15}
\]
This proves both split local conditions when there are two,
and the full unique inert local condition otherwise.
No condition at a second split place has been dropped.

\proofstep{4. The original local invariant quotients and uniformizer.}
For each $v\mid\ell$, a Frobenius lift has scalar on $D$
whose reduction is (MC7.15), even if $\kappa$ is ramified.
Its scalar minus one is an $O$-unit, so $D^{G_{K_v}}=0$.
At any finite $v\nmid\ell$, both $c_\ell$ and $\kappa$
are unramified.  The finite inertia image of $\chi^j\nu^a$
has order dividing three.  If it is nontrivial, some scalar
minus one is an $O$-unit and $D^{I_v}=0$; otherwise
$D^{I_v}=D$ is $\varpi$-divisible.  This includes the
possibly ramified places over $3,7,269$.  Consequently
\[
 D^{G_{K_v}}/\varpi D^{G_{K_v}}=0\ (v\mid\ell),
 \qquad D^{I_v}/\varpi D^{I_v}=0\ (v\nmid\ell).
 \tag{MC7.16}
\]
For every one of these groups $G$, the original coefficient
sequence gives
\[
 D^G/\varpi D^G\longrightarrow H^1(G,D[\varpi])
            \longrightarrow H^1(G,D)[\varpi]\longrightarrow0.
 \tag{MC7.17}
\]
Lift a global strict class killed by $\varpi$ through the
global surjection in (MC7.17).  Its local residual classes
map to zero, and (MC7.16) shows that they are zero already.
The lift is full strict residually and vanishes by Step 3.
Thus the original strict group has no $\varpi$-torsion.
Every continuous cocycle into the discrete $E/O$ has finite
image, since its domain is compact.  It is therefore killed
by some power of $\varpi$.  Any nonzero strict class would
produce nonzero $\varpi$-torsion, a contradiction.  This
proves the second assertion of (MC7.4), with no coefficient
extension or auxiliary local weakening.

\proofstep{5. Actual growing hosts, exact counts and retained sums.}
The product decomposition (FC.1) contains independently the
Kummer $C_3$ and the rational cubic subfactor of its conductor-seven
$C_6$.  Both are present for every actual $h\geq1$.
Thus the four displayed products are distinct actual original
characters.  Their values lie in the original completion of
$\mathbb Z[\zeta_{3\cdot2^h}]$, since only $\mu_3$ is needed.
No new finite cyclotomic character at the coefficient prime is
inserted into the source host.

Set $\kappa=1$ and identify the original module in (FC.2)
with $(E_\ell/O_\ell)(\chi^j\xi_7^{2a}c_\ell^{-e})$.
Inflation from the actual $G_{K,\Sigma_\ell}$ to $G_K$ is
injective.  Its strict classes are inertia-zero outside
$\Sigma_\ell$ and retain exactly the same full local kernels
inside it, so they inject into the zero group in (MC7.4).
The original $O_\ell$, its uniformizer and normalizing rank
$b_\ell$ are unchanged.

For each eligible prime the number of positive weights
$e<T$ divisible by $\ell-1$ is
$\lfloor(T-1)/(\ell-1)\rfloor$, since $T$ is an integer.
It is nonzero only when $\ell\leq T$.  Multiplying by the
four distinct labels proves (MC7.6).  The primes five and
eleven are inert in $K$ and not congruent to $\pm1$ modulo
seven, so they have $c=0,s=1$.  At thirteen,
$13\equiv1\pmod3$, $13\equiv-1\pmod7$ and
$269^4\equiv9\ne1\pmod{13}$, giving $c=1,s=0$.
These are the three stated examples.

These labels have nontrivial cubic Kummer factor and a
nontrivial rational cubic seven-factor.  In particular they
are distinct from the pure cubic labels of AC-G203 and the
quadratic-twisted labels of AC-G207, whose rational seven-factor
is trivial.  Their contribution to (CP.3) is a full original
$K$-strict summand, and has just been proved zero.  Deleting
these summands from any $J$ gives (MC7.7), with no new
comparison cost or separately chosen window.
\end{proof}

\paragraph{Finite-cover dictionary and source scope.}
For every $n\geq1$, the original $D[\varpi^n]$ defines a
finite locally constant etale sheaf on
$\operatorname{Spec}\mathcal O_K[1/(3\cdot7\cdot269\cdot\ell)]$.
Its residual action is trivialized by the stated degree-nine
$C_3\times C_3$ cover.  The full strict kernels are the
original completion and inertia restrictions on this sheaf.
This is the retained SGA~1, Expose V, sections 4 and 7
finite-cover/Artin--Tate interpretation; no new archival read
or motivic realization is asserted.
The root proposed the split noncube branch and owns the exact
class execution.  The author supplies the complementary local
case, the resulting full prime criterion, and the complete
degree-two transposition, residual-cover and original-coefficient proof,
using the independently checked G203/G207 mechanism.
The complete applied Sharifi class-field/norm passage and
documented PARI flag-one scope are retained explicitly.
The all-prime remainder, $\mathcal R_{\rm cub}$ and H04.4
remain separate main gates.
